\documentclass[12pt]{article}
\usepackage[spanish]{babel}
\usepackage[utf8]{inputenc}
\usepackage[T1]{fontenc}
\usepackage{graphicx}
\usepackage{geometry}
\usepackage{setspace}
\usepackage[hidelinks]{hyperref}
\usepackage{float}
\usepackage{titlesec}
\usepackage{listings}
\usepackage{xcolor}
\usepackage{enumitem}
\usepackage{amsmath}
\usepackage{booktabs}
\usepackage{longtable}
\usepackage{array}
\usepackage{multirow}

\geometry{a4paper, margin=2.5cm}

% --- Configuracion para codigo ---
\definecolor{codegreen}{rgb}{0,0.6,0}
\definecolor{codegray}{rgb}{0.5,0.5,0.5}
\definecolor{backcolour}{rgb}{0.95,0.95,0.92}

\lstset{
    backgroundcolor=\color{backcolour},   
    commentstyle=\color{codegreen},
    keywordstyle=\color{blue},
    numberstyle=\tiny\color{codegray},
    stringstyle=\color{orange},
    basicstyle=\ttfamily\footnotesize,
    breaklines=true,
    captionpos=b,
    numbers=left,
}

\begin{document}
% --- PORTADA ---
\begin{titlepage}
\begin{center}
    \includegraphics[width=0.25\textwidth]{logo-uteq.png} \\[1cm]
    {\large \textbf{UNIVERSIDAD TÉCNICA ESTATAL DE QUEVEDO}} \\[0.5cm]
    {\large \textbf{FACULTAD DE CIENCIAS DE LA INGENIERÍA}} \\[0.5cm]
    {\large \textbf{CARRERA INGENIERÍA DE SOFTWARE (REDISEÑO)}} \\[0.5cm]
    
    \textbf{TEMA:}\\[0.5cm]
    {\Large \textbf{Simulador de Comportamiento Vial con Inteligencia Artificial (SBVIA)}} \\[0.5cm]
    {\large Entrega 1A - Planificación y Diseño} \\[0.5cm]
    
    \textbf{GRUPO:}\\[0.5cm]
    \textbf{CRUZ PEREZ JUSTYN KEITH}\\
    \textbf{UMAGINGA AREVALO JEFFERSON MANUEL}\\
    \textbf{ZAMORA BUMBILA DIEGO ALEXANDER} \\[1cm]
    
    \textbf{DOCENTE:}\\[0.5cm]
    ING. GLEISTON CICERON GUERRERO ULLOA\\[1cm]
    
    \textbf{MATERIA:}\\[0.5cm]
    Aplicaciones Web - A - [5to Nivel] \\[1cm]
    
    \textbf{URL DEL REPOSITORIO:}\\[0.5cm]
    \url{https://github.com/keithdrox/SBVIA-AppWeb} \\[1cm]

    \vfill
    Quevedo - Ecuador \\
    2026
\end{center}
\end{titlepage}

\newpage
\tableofcontents
\newpage

\section{Descripción del sistema web propuesto}
El comportamiento vial de los conductores en formación presenta errores recurrentes que los sistemas de enseñanza tradicionales no logran prevenir de forma sistemática. Acciones como el exceso de velocidad, el incumplimiento de señales de tránsito y la conducción distraída concentran una proporción significativa de los incidentes en entornos urbanos del Ecuador. 

El sistema \textbf{Simulador de Comportamiento Vial con Inteligencia Artificial (SBVIA)} aborda este problema proporcionando un entorno de entrenamiento y evaluación. Los usuarios principales del sistema son: \textbf{Conductores en formación} (que interactúan con el simulador), \textbf{Instructores de conducción} (que supervisan el desempeño), \textbf{Administradores} (que configuran escenarios) y \textbf{Autoridades de tránsito}.

Se ha optado por implementar SBVIA como una \textbf{aplicación web} porque facilita el acceso ubicuo desde cualquier dispositivo con navegador, eliminando la necesidad de instalaciones complejas de software de simulación en los equipos de los sindicatos de choferes. Además, permite la centralización de las métricas de evaluación en la nube. 

En esta entrega, el alcance contempla el núcleo funcional desplegable: registro y autenticación de usuarios, selección y configuración de escenarios de conducción, registro de decisiones y generación de retroalimentación automatizada y métricas de desempeño.

\section{Revisión y actualización de requisitos del semestre anterior}
\subsection{Resumen del trabajo previo}
En el semestre anterior (Ingeniería de Requerimientos), el equipo aplicó el marco de trabajo ágil Scrum y el enfoque GORE (Goal-Oriented Requirements Engineering). Se emplearon técnicas de elicitación como entrevistas estructuradas, análisis documental y observación del dominio en el Sindicato de Choferes del Cantón Valencia. Se documentaron 11 Requisitos Funcionales (RF) y 6 Requisitos No Funcionales (RNF), priorizados mediante la técnica MoSCoW, junto con el modelado UML completo y validación con los stakeholders.

\subsection{Tabla de adaptación de requisitos}
\begin{table}[H]
\centering
\resizebox{\textwidth}{!}{
\begin{tabular}{|p{1.5cm}|p{6cm}|p{2cm}|p{6cm}|}
\hline
\textbf{Cód.} & \textbf{Requisito original (semestre ant.)} & \textbf{Estado} & \textbf{Adaptación para la aplicación web} \\ \hline
RF-01 & Autenticación de usuarios mediante credenciales válidas. & Mantiene & Se implementa tal cual como endpoint REST y formulario web. \\ \hline
RF-02 & Seleccionar escenarios de simulación con niveles de dificultad. & Adapta & Interfaz web para elegir escenarios; carga dinámica de recursos. \\ \hline
RF-03 & Simular señales, peatones y vehículos en entornos urbanos. & Adapta & Renderizado web (WebGL/Canvas) del entorno de simulación. \\ \hline
RF-04 & Registrar decisiones en tiempo real. & Adapta & Comunicación asíncrona (WebSockets o Fetch) con el backend. \\ \hline
RF-05 & Analizar comportamiento usando IA. & Mantiene & Procesamiento en servidor (backend) y respuesta vía API. \\ \hline
RF-06 & Generar métricas cuantificables. & Mantiene & Visualización en dashboards web interactivos (ej. Chart.js). \\ \hline
RF-07 & Proporcionar retroalimentación educativa al finalizar. & Adapta & Generación de reporte web dinámico e interactivo. \\ \hline
RF-08 & Almacenar historial de prácticas. & Mantiene & Almacenamiento en base de datos centralizada. \\ \hline
\end{tabular}}
\end{table}

\subsection{Nuevos requisitos específicos del contexto web}
\begin{itemize}
    \item \textbf{RNF-WEB-01:} La aplicación debe ser accesible desde navegadores modernos (Chrome 120+, Firefox 120+, Edge 120+).
    \item \textbf{RNF-WEB-02:} El tiempo de respuesta inicial del dashboard no debe superar los 3 segundos.
    \item \textbf{RNF-WEB-03:} La interfaz de reportes debe ser responsiva en pantallas de 320px a 1440px.
    \item \textbf{RNF-WEB-04:} Todas las comunicaciones con la API de IA deben realizarse sobre HTTPS.
    \item \textbf{RF-WEB-01:} El sistema debe gestionar sesiones de usuario de forma segura utilizando JWT o sesiones tradicionales de backend.
\end{itemize}

\section{Lista consolidada de requisitos funcionales}
\begin{longtable}{|p{1.5cm}|p{5cm}|p{2cm}|p{6cm}|}
\hline
\textbf{ID} & \textbf{Descripción} & \textbf{Prioridad} & \textbf{Criterio de aceptación} \\ \hline
\endfirsthead
\hline
\textbf{ID} & \textbf{Descripción} & \textbf{Prioridad} & \textbf{Criterio de aceptación} \\ \hline
\endhead
RF-01 & El sistema deberá autenticar a los usuarios mediante credenciales válidas. & Alta & El usuario ingresa credenciales válidas y es redirigido al dashboard. \\ \hline
RF-02 & El sistema deberá permitir al usuario seleccionar escenarios de simulación (básico, intermedio, avanzado). & Alta & El usuario puede iniciar un escenario y los parámetros se cargan en < 2s. \\ \hline
RF-03 & El sistema deberá simular señales, peatones y vehículos en el entorno virtual en un máximo de 500ms tras iniciar. & Alta & Los obstáculos y señales son visibles y ocupan sus coordenadas 3D en < 500ms. \\ \hline
RF-04 & El sistema deberá registrar en tiempo real las decisiones del conductor. & Alta & Las acciones (frenado, giro) se envían al backend con latencia < 100ms. \\ \hline
RF-05 & El sistema deberá analizar el comportamiento de conducción mediante IA, retornando una clasificación de riesgo (Bajo, Medio, Alto) en máximo 2s. & Alta & El sistema evalúa las acciones registradas y retorna la clasificación de riesgo en < 2s. \\ \hline
RF-06 & El sistema deberá calcular el puntaje de la sesión. & Media & Se calcula el puntaje final tras concluir la sesión. \\ \hline
RF-06b & El sistema deberá contabilizar las infracciones cometidas en la sesión. & Media & Se registran las infracciones individuales durante la sesión. \\ \hline
RF-06c & El sistema deberá calcular el tiempo de reacción del conductor. & Media & Se calcula el tiempo de reacción promedio al finalizar la sesión. \\ \hline
RF-07 & El sistema deberá generar un reporte educativo listando cada error y su respectiva regla de tránsito en menos de 3s. & Alta & Se muestra el reporte con cada error y regla infringida en < 3s. \\ \hline
RF-08 & El sistema deberá almacenar el historial de prácticas de cada usuario. & Media & El usuario puede ver al menos 5 sesiones pasadas en su perfil. \\ \hline
RF-09 & El sistema deberá proveer un resumen de métricas agregadas a los instructores. & Media & El instructor visualiza métricas de todos sus alumnos. \\ \hline
RF-10 & El sistema deberá permitir a los administradores crear y editar reglas de tránsito. & Baja & El admin puede gestionar reglas desde un panel de administración. \\ \hline
\end{longtable}

\section{Lista consolidada de requisitos no funcionales}
\begin{table}[H]
\centering
\begin{tabular}{|p{2cm}|p{2.5cm}|p{5cm}|p{5cm}|}
\hline
\textbf{ID} & \textbf{Categoría} & \textbf{Descripción} & \textbf{Métrica de verificación} \\ \hline
RNF-01 & Rendimiento & Latencia de registro de decisiones inferior a 100ms. & Medición en Network Tab (Chrome DevTools) P95 < 100ms. \\ \hline
RNF-02 & Seguridad & Datos de usuarios protegidos y contraseñas hasheadas. & Revisión en BD de la columna contrasena\_hash en Argon2id o BCrypt. \\ \hline
RNF-03 & Usabilidad & Interfaz responsiva adaptable a múltiples pantallas. & 100\% de visibilidad sin scroll horizontal en resoluciones de 320px a 1440px. \\ \hline
RNF-04 & Disponibilidad & El sistema debe estar disponible para demostraciones. & Uptime del 99\% durante las semanas 16 y 17 de clases. \\ \hline
RNF-05 & Mantenibilidad & El código backend debe seguir principios de código limpio. & Análisis estático sin violaciones críticas. \\ \hline
RNF-06 & Escalabilidad & El sistema debe soportar acceso concurrente de alumnos. & Pruebas k6 con 50 VUs manteniendo tiempos de respuesta < 200ms. \\ \hline
RNF-07 & Portabilidad & El sistema debe ser desplegable en múltiples entornos. & Despliegue mediante Docker Compose en cualquier SO con Docker. \\ \hline
\end{tabular}
\end{table}

\section{Arquitectura de la aplicación web}
\subsection{Diagramas de arquitectura}


\begin{figure}[H]
    \centering
    \includegraphics[width=0.8\textwidth]{C4 Nivel 1.jpg}
    \caption{Diagrama C4 Nivel 1 (Contexto) - Sistema SBVIA interactuando con actores y sistemas externos.}
\end{figure}

El diagrama muestra cómo los usuarios y administradores interactúan con el Simulador de Comportamiento Vial con Inteligencia Artificial. El sistema se integra con una API de IA para evaluar el comportamiento de conducción y con un servicio de correo electrónico para enviar notificaciones y recuperación de contraseñas.


\begin{figure}[H]
    \centering
    \includegraphics[width=0.8\textwidth]{C4 Nivel 2.jpg}
    \caption{Diagrama C4 Nivel 2 (Contenedores) - Aplicación Web, API Backend y Base de Datos.}
\end{figure}
Describe la estructura interna del sistema compuesta por un Frontend desarrollado con HTML5, CSS3, JavaScript y Bootstrap 5; un Backend implementado con PHP 8.2 y Laravel 11; y una base de datos PostgreSQL 17. También se muestran las integraciones externas con la API de Inteligencia Artificial y el servicio de correo electrónico.

\subsection{Descripción de capas y tecnologías}
\begin{table}[H]
\centering
\begin{tabular}{|p{3cm}|p{4cm}|p{7cm}|}
\hline
\textbf{Capa} & \textbf{Tecnología elegida} & \textbf{Justificación} \\ \hline
Frontend & HTML5 + CSS3 + JS ES6+ & Permite crear interfaces dinámicas e interactivas. \\ \hline
Framework CSS & Bootstrap 5 & Facilita el desarrollo responsivo y mejora la experiencia de usuario. \\ \hline
Backend & PHP 8.2 & Amplio soporte y compatibilidad con aplicaciones web modernas. \\ \hline
Framework Backend & Laravel 11 & Implementa MVC, ORM Eloquent, autenticación y seguridad integrada. \\ \hline
Base de Datos & PostgreSQL 17 & Motor robusto, seguro, escalable y con excelente manejo de consultas. \\ \hline
Servidor Web & Apache & Fácil configuración e integración con PHP. \\ \hline
Control de Versiones & Git + GitHub & Gestión colaborativa del proyecto. \\ \hline
Despliegue & XAMPP / Docker & Facilita el desarrollo y pruebas del sistema. \\ \hline
\end{tabular}
\end{table}

\subsection{ADR-001: Decisión de tecnología principal}
\begin{table}[H]
\centering
\begin{tabular}{|p{3cm}|p{11cm}|}
\hline
\textbf{Título} & \textbf{ADR-001: Selección del Lenguaje de Programación del Servidor} \\ \hline
\textbf{Estado} & Aceptado \\ \hline
\textbf{Contexto} & El equipo necesita seleccionar la tecnología del lado del servidor para desarrollar el sistema "Simulador de Comportamiento Vial con Inteligencia Artificial". La decisión debe considerar el tiempo disponible para el desarrollo, el conocimiento previo de los integrantes, la integración con PostgreSQL, la compatibilidad con APIs de Inteligencia Artificial y la facilidad de despliegue en servidores web. \\ \hline
\textbf{Opciones consideradas} & Opción A: PHP 8.2 con Laravel 11. Opción B: ASP .NET Core 8. Opción C: PHP 8.2 sin framework. \\ \hline
\textbf{Decisión} & Se selecciona \textbf{PHP 8.2 con Laravel 11} debido a que proporciona una arquitectura MVC robusta, integración sencilla con PostgreSQL, herramientas de autenticación y seguridad integradas, facilidad para consumir APIs externas y una curva de aprendizaje adecuada para el equipo. \\ \hline
\textbf{Consecuencias positivas} & \begin{itemize}[leftmargin=*,nosep]
    \item Desarrollo más rápido mediante componentes predefinidos.
    \item Arquitectura MVC organizada.
    \item Fácil integración con PostgreSQL.
    \item Mayor seguridad frente a ataques comunes.
    \item Amplia documentación y comunidad de soporte.
    \item Facilidad de integración con servicios de Inteligencia Artificial.
\end{itemize} \\ \hline
\textbf{Consecuencias negativas} & \begin{itemize}[leftmargin=*,nosep]
    \item Dependencia del framework Laravel.
    \item Necesidad de capacitación inicial para algunos integrantes.
    \item Mayor consumo de recursos que PHP puro.
    \item Posibles cambios de compatibilidad en futuras versiones del framework.
\end{itemize} \\ \hline
\end{tabular}
\end{table}

\section{Modelo de base de datos}
\subsection{Modelo Entidad-Relación (MER)}


\begin{figure}[H]
    \centering
    \includegraphics[width=0.8\textwidth]{waos.jpg}
    \caption{Diagrama Entidad-Relación del Sistema SBVIA.}
\end{figure}

\subsection{Diccionario de datos}

\textbf{Tabla: Rol}
\begin{table}[H]
\centering
\begin{tabular}{|p{3cm}|p{2cm}|p{1cm}|p{2cm}|p{5cm}|}
\hline
\textbf{Atributo} & \textbf{Tipo} & \textbf{Nulo} & \textbf{Restricción} & \textbf{Descripción} \\ \hline
id\_Rol & INTEGER & No & PK, AUTO & Identificador único del rol \\ \hline
nombre & VARCHAR & No & NOT NULL & Nombre del rol \\ \hline
descripcion & VARCHAR & No & NOT NULL & Descripción del rol \\ \hline
\end{tabular}
\end{table}

\textbf{Tabla: Usuario}
\begin{table}[H]
\centering
\begin{tabular}{|p{3cm}|p{2cm}|p{1cm}|p{2cm}|p{5cm}|}
\hline
\textbf{Atributo} & \textbf{Tipo} & \textbf{Nulo} & \textbf{Restricción} & \textbf{Descripción} \\ \hline
id\_Usuario & INTEGER & No & PK, AUTO & Identificador único del usuario \\ \hline
nombre & VARCHAR & No & NOT NULL & Nombre del usuario \\ \hline
apellido & VARCHAR & No & NOT NULL & Apellido del usuario \\ \hline
fecha\_registro & DATE & No & NOT NULL & Fecha de registro \\ \hline
estado & VARCHAR & No & NOT NULL & Estado (Activo/Inactivo) \\ \hline
id\_Rol & INTEGER & No & FK & Rol asignado \\ \hline
\end{tabular}
\end{table}

\textbf{Tabla: Escenario}
\begin{table}[H]
\centering
\begin{tabular}{|p{3cm}|p{2cm}|p{1cm}|p{2cm}|p{5cm}|}
\hline
\textbf{Atributo} & \textbf{Tipo} & \textbf{Nulo} & \textbf{Restricción} & \textbf{Descripción} \\ \hline
id\_Escenario & INTEGER & No & PK, AUTO & Identificador único del escenario \\ \hline
nombre & VARCHAR & No & NOT NULL & Nombre del escenario vial \\ \hline
descripcion & VARCHAR & Sí & - & Descripción del escenario \\ \hline
tipo\_via & VARCHAR & No & NOT NULL & Tipo de vía (urbana, rural, etc.) \\ \hline
nivel\_dificultad & INTEGER & No & NOT NULL & Nivel de dificultad \\ \hline
clima & VARCHAR & No & NOT NULL & Condición climática \\ \hline
densidad\_trafico & VARCHAR & No & NOT NULL & Nivel de tráfico \\ \hline
\end{tabular}
\end{table}

\textbf{Tabla: Simulacion}
\begin{table}[H]
\centering
\begin{tabular}{|p{3cm}|p{2cm}|p{1cm}|p{2cm}|p{5cm}|}
\hline
\textbf{Atributo} & \textbf{Tipo} & \textbf{Nulo} & \textbf{Restricción} & \textbf{Descripción} \\ \hline
id\_Simulacion & INTEGER & No & PK, AUTO & Identificador de la simulación \\ \hline
fecha\_inicio & DATE & No & NOT NULL & Fecha de inicio \\ \hline
fecha\_fin & DATE & No & NOT NULL & Fecha de finalización \\ \hline
estado & VARCHAR & No & NOT NULL & Estado de la simulación \\ \hline
puntaje\_final & DECIMAL & No & NOT NULL & Puntaje obtenido \\ \hline
id\_Usuario & INTEGER & No & FK & Usuario que realiza la simulación \\ \hline
id\_Escenario & INTEGER & No & FK & Escenario utilizado \\ \hline
\end{tabular}
\end{table}

\textbf{Tabla: EventoVial}
\begin{table}[H]
\centering
\begin{tabular}{|p{3cm}|p{2cm}|p{1cm}|p{2cm}|p{5cm}|}
\hline
\textbf{Atributo} & \textbf{Tipo} & \textbf{Nulo} & \textbf{Restricción} & \textbf{Descripción} \\ \hline
id\_EventoVial & INTEGER & No & PK, AUTO & Identificador del evento \\ \hline
nombre & VARCHAR & No & NOT NULL & Nombre del evento \\ \hline
descripcion & VARCHAR & Sí & - & Descripción del evento \\ \hline
tipo\_evento & VARCHAR & No & NOT NULL & Tipo de evento vial \\ \hline
nivel\_riesgo & INTEGER & No & NOT NULL & Nivel de riesgo \\ \hline
id\_Simulacion & INTEGER & No & FK & Simulación asociada \\ \hline
\end{tabular}
\end{table}

\textbf{Tabla: Decision}
\begin{table}[H]
\centering
\begin{tabular}{|p{3cm}|p{2cm}|p{1cm}|p{2cm}|p{5cm}|}
\hline
\textbf{Atributo} & \textbf{Tipo} & \textbf{Nulo} & \textbf{Restricción} & \textbf{Descripción} \\ \hline
id\_Decision & INTEGER & No & PK, AUTO & Identificador de la decisión \\ \hline
accion\_realizada & VARCHAR & No & NOT NULL & Acción tomada por el usuario \\ \hline
resultado & VARCHAR & No & NOT NULL & Resultado obtenido \\ \hline
momento & TIMESTAMP & No & NOT NULL & Momento de la acción \\ \hline
id\_Simulacion & INTEGER & No & FK & Simulación relacionada \\ \hline
\end{tabular}
\end{table}

\textbf{Tabla: ReglaTransito}
\begin{table}[H]
\centering
\begin{tabular}{|p{3cm}|p{2cm}|p{1cm}|p{2cm}|p{5cm}|}
\hline
\textbf{Atributo} & \textbf{Tipo} & \textbf{Nulo} & \textbf{Restricción} & \textbf{Descripción} \\ \hline
id\_ReglaTransito & INTEGER & No & PK, AUTO & Identificador de la regla \\ \hline
nombre & VARCHAR & No & NOT NULL & Nombre de la norma \\ \hline
descripcion & VARCHAR & Sí & - & Descripción de la regla \\ \hline
categoria & VARCHAR & No & NOT NULL & Categoría de la norma \\ \hline
id\_Escenario & INTEGER & No & FK & Escenario donde aplica \\ \hline
\end{tabular}
\end{table}

\textbf{Tabla: Infraccion}
\begin{table}[H]
\centering
\begin{tabular}{|p{3cm}|p{2cm}|p{1cm}|p{2cm}|p{5cm}|}
\hline
\textbf{Atributo} & \textbf{Tipo} & \textbf{Nulo} & \textbf{Restricción} & \textbf{Descripción} \\ \hline
id\_Infraccion & INTEGER & No & PK, AUTO & Identificador de la infracción \\ \hline
nombre & VARCHAR & No & NOT NULL & Nombre de la infracción \\ \hline
descripcion & VARCHAR & Sí & - & Descripción de la infracción \\ \hline
gravedad & VARCHAR & No & NOT NULL & Nivel de gravedad \\ \hline
penalizacion & DECIMAL & No & NOT NULL & Penalización aplicada \\ \hline
id\_Simulacion & INTEGER & No & FK & Simulación asociada \\ \hline
\end{tabular}
\end{table}

\textbf{Tabla: ComportamientoVial}
\begin{table}[H]
\centering
\begin{tabular}{|p{3cm}|p{2cm}|p{1cm}|p{2cm}|p{5cm}|}
\hline
\textbf{Atributo} & \textbf{Tipo} & \textbf{Nulo} & \textbf{Restricción} & \textbf{Descripción} \\ \hline
id\_ComportamientoVial & INTEGER & No & PK, AUTO & Identificador del comportamiento \\ \hline
clasificacion & VARCHAR & No & NOT NULL & Clasificación \\ \hline
nivel\_riesgo & INTEGER & No & NOT NULL & Nivel de riesgo detectado \\ \hline
observaciones & VARCHAR & Sí & - & Observaciones de la IA \\ \hline
id\_Decision & INTEGER & No & FK & Decisión evaluada \\ \hline
\end{tabular}
\end{table}

\textbf{Tabla: Retroalimentacion}
\begin{table}[H]
\centering
\begin{tabular}{|p{3cm}|p{2cm}|p{1cm}|p{2cm}|p{5cm}|}
\hline
\textbf{Atributo} & \textbf{Tipo} & \textbf{Nulo} & \textbf{Restricción} & \textbf{Descripción} \\ \hline
id\_Retroalimentacion & INTEGER & No & PK, AUTO & Identificador \\ \hline
comentario & VARCHAR & No & NOT NULL & Comentario generado por IA \\ \hline
recomendacion & VARCHAR & Sí & - & Recomendación \\ \hline
fecha\_generacion & DATE & No & NOT NULL & Fecha de generación \\ \hline
id\_Simulacion & INTEGER & No & FK & Simulación relacionada \\ \hline
id\_ComportamientoVial & INTEGER & No & FK & Comportamiento evaluado \\ \hline
\end{tabular}
\end{table}

\textbf{Tabla: MetricaDesempeno}
\begin{table}[H]
\centering
\begin{tabular}{|p{3cm}|p{2cm}|p{1cm}|p{2cm}|p{5cm}|}
\hline
\textbf{Atributo} & \textbf{Tipo} & \textbf{Nulo} & \textbf{Restricción} & \textbf{Descripción} \\ \hline
id\_MetricaDesempeno & INTEGER & No & PK, AUTO & Identificador de métrica \\ \hline
puntaje & DECIMAL & No & NOT NULL & Puntaje obtenido \\ \hline
tiempo\_reaccion & DECIMAL & No & NOT NULL & Tiempo promedio de reacción \\ \hline
errores & INTEGER & No & NOT NULL & Errores cometidos \\ \hline
aciertos & INTEGER & No & NOT NULL & Aciertos \\ \hline
nivel\_desempeno & VARCHAR & No & NOT NULL & Nivel de desempeño \\ \hline
id\_Simulacion & INTEGER & No & FK & Simulación evaluada \\ \hline
\end{tabular}
\end{table}

\textbf{Tabla: Reporte}
\begin{table}[H]
\centering
\begin{tabular}{|p{3cm}|p{2cm}|p{1cm}|p{2cm}|p{5cm}|}
\hline
\textbf{Atributo} & \textbf{Tipo} & \textbf{Nulo} & \textbf{Restricción} & \textbf{Descripción} \\ \hline
id\_Reporte & INTEGER & No & PK, AUTO & Identificador del reporte \\ \hline
tipo\_reporte & VARCHAR & No & NOT NULL & Tipo de reporte \\ \hline
fecha\_generacion & DATE & No & NOT NULL & Fecha de generación \\ \hline
observaciones & VARCHAR & No & NOT NULL & Observaciones generales \\ \hline
id\_Simulacion & INTEGER & No & FK & Simulación asociada \\ \hline
\end{tabular}
\end{table}

\textbf{Tabla: Recompensa}
\begin{table}[H]
\centering
\begin{tabular}{|p{3cm}|p{2cm}|p{1cm}|p{2cm}|p{5cm}|}
\hline
\textbf{Atributo} & \textbf{Tipo} & \textbf{Nulo} & \textbf{Restricción} & \textbf{Descripción} \\ \hline
id\_Recompensa & INTEGER & No & PK, AUTO & Identificador de la recompensa \\ \hline
nombre & VARCHAR & No & NOT NULL & Nombre \\ \hline
descripcion & VARCHAR & No & NOT NULL & Descripción \\ \hline
nivel & INTEGER & No & NOT NULL & Nivel requerido \\ \hline
puntos\_requeridos & INTEGER & No & NOT NULL & Puntos necesarios \\ \hline
id\_Usuario & INTEGER & No & FK & Usuario que recibe recompensa \\ \hline
\end{tabular}
\end{table}

\subsection{Script DDL de creación}
\begin{lstlisting}[language=SQL]
CREATE TABLE "Rol" ( 
  "id_Rol" INTEGER GENERATED BY DEFAULT AS IDENTITY PRIMARY KEY NOT NULL, 
  "nombre" varchar NOT NULL, 
  "descripcion" varchar NOT NULL 
); 

CREATE TABLE "Usuario" ( 
  "id_Usuario" INTEGER GENERATED BY DEFAULT AS IDENTITY PRIMARY KEY NOT NULL, 
  "nombre" varchar NOT NULL, 
  "apellido" varchar NOT NULL, 
  "fecha_registro" date NOT NULL, 
  "estado" varchar NOT NULL, 
  "id_Rol" integer NOT NULL 
); 

CREATE TABLE "Escenario" ( 
  "id_Escenario" INTEGER GENERATED BY DEFAULT AS IDENTITY PRIMARY KEY NOT NULL, 
  "nombre" varchar NOT NULL, 
  "descripcion" varchar, 
  "tipo_via" varchar NOT NULL, 
  "nivel_dificultad" integer NOT NULL, 
  "clima" varchar NOT NULL, 
  "densidad_trafico" varchar NOT NULL 
); 

CREATE TABLE "Simulacion" ( 
  "id_Simulacion" INTEGER GENERATED BY DEFAULT AS IDENTITY PRIMARY KEY NOT NULL, 
  "fecha_inicio" date NOT NULL, 
  "fecha_fin" date NOT NULL, 
  "estado" varchar NOT NULL, 
  "puntaje_final" decimal NOT NULL, 
  "id_Usuario" integer NOT NULL, 
  "id_Escenario" integer NOT NULL 
); 

CREATE TABLE "EventoVial" ( 
  "id_EventoVial" INTEGER GENERATED BY DEFAULT AS IDENTITY PRIMARY KEY NOT NULL, 
  "nombre" varchar NOT NULL, 
  "descripcion" varchar, 
  "tipo_evento" varchar NOT NULL, 
  "nivel_riesgo" integer NOT NULL, 
  "id_Simulacion" integer NOT NULL 
); 

CREATE TABLE "Decision" ( 
  "id_Decision" INTEGER GENERATED BY DEFAULT AS IDENTITY PRIMARY KEY NOT NULL, 
  "accion_realizada" varchar NOT NULL, 
  "resultado" varchar NOT NULL, 
  "momento" date NOT NULL, 
  "id_Simulacion" integer NOT NULL 
); 

CREATE TABLE "ReglaTransito" ( 
  "id_ReglaTransito" INTEGER GENERATED BY DEFAULT AS IDENTITY PRIMARY KEY NOT NULL, 
  "nombre" varchar NOT NULL, 
  "descripcion" varchar, 
  "categoria" varchar NOT NULL, 
  "id_Escenario" integer NOT NULL 
); 

CREATE TABLE "Infraccion" ( 
  "id_Infraccion" INTEGER GENERATED BY DEFAULT AS IDENTITY PRIMARY KEY NOT NULL, 
  "nombre" varchar NOT NULL, 
  "descripcion" varchar, 
  "gravedad" varchar NOT NULL, 
  "penalizacion" decimal NOT NULL, 
  "id_Simulacion" integer NOT NULL 
); 

CREATE TABLE "ComportamientoVial" ( 
  "id_ComportamientoVial" INTEGER GENERATED BY DEFAULT AS IDENTITY PRIMARY KEY NOT NULL, 
  "clasificacion" varchar NOT NULL, 
  "nivel_riesgo" integer NOT NULL, 
  "observaciones" varchar, 
  "id_Decision" integer NOT NULL 
); 

CREATE TABLE "Retroalimentacion" ( 
  "id_Retroalimentacion" INTEGER GENERATED BY DEFAULT AS IDENTITY PRIMARY KEY NOT NULL, 
  "comentario" varchar NOT NULL, 
  "recomendacion" varchar, 
  "fecha_generacion" date NOT NULL, 
  "id_Simulacion" integer NOT NULL, 
  "id_ComportamientoVial" integer NOT NULL 
); 

CREATE TABLE "MetricaDesempeno" ( 
  "id_MetricaDesempeno" INTEGER GENERATED BY DEFAULT AS IDENTITY PRIMARY KEY NOT NULL, 
  "puntaje" decimal NOT NULL, 
  "tiempo_reaccion" decimal NOT NULL, 
  "errores" integer NOT NULL, 
  "aciertos" integer NOT NULL, 
  "nivel_desempeno" varchar NOT NULL, 
  "id_Simulacion" integer NOT NULL 
); 

CREATE TABLE "Reporte" ( 
  "id_Reporte" INTEGER GENERATED BY DEFAULT AS IDENTITY PRIMARY KEY NOT NULL, 
  "tipo_reporte" varchar NOT NULL, 
  "fecha_generacion" date NOT NULL, 
  "observaciones" varchar NOT NULL, 
  "id_Simulacion" integer NOT NULL 
); 

CREATE TABLE "Recompensa" ( 
  "id_Recompensa" INTEGER GENERATED BY DEFAULT AS IDENTITY PRIMARY KEY NOT NULL, 
  "nombre" varchar NOT NULL, 
  "descripcion" varchar NOT NULL, 
  "nivel" integer NOT NULL, 
  "puntos_requeridos" integer NOT NULL, 
  "id_Usuario" integer NOT NULL 
); 

ALTER TABLE "Simulacion" ADD CONSTRAINT "fk_Simulacion_id_Escenario_Escenario" FOREIGN KEY ("id_Escenario") REFERENCES "Escenario" ("id_Escenario") ON DELETE NO ACTION ON UPDATE NO ACTION DEFERRABLE INITIALLY IMMEDIATE; 
ALTER TABLE "Usuario" ADD CONSTRAINT "fk_Usuario_id_Rol_Rol" FOREIGN KEY ("id_Rol") REFERENCES "Rol" ("id_Rol") ON DELETE NO ACTION ON UPDATE NO ACTION DEFERRABLE INITIALLY IMMEDIATE; 
ALTER TABLE "Simulacion" ADD CONSTRAINT "fk_Simulacion_id_Usuario_Usuario" FOREIGN KEY ("id_Usuario") REFERENCES "Usuario" ("id_Usuario") ON DELETE NO ACTION ON UPDATE NO ACTION DEFERRABLE INITIALLY IMMEDIATE; 
ALTER TABLE "Decision" ADD CONSTRAINT "fk_Decision_id_Simulacion_Simulacion" FOREIGN KEY ("id_Simulacion") REFERENCES "Simulacion" ("id_Simulacion") ON DELETE NO ACTION ON UPDATE NO ACTION DEFERRABLE INITIALLY IMMEDIATE; 
ALTER TABLE "ComportamientoVial" ADD CONSTRAINT "fk_ComportamientoVial_id_Decision_Decision" FOREIGN KEY ("id_Decision") REFERENCES "Decision" ("id_Decision") ON DELETE NO ACTION ON UPDATE NO ACTION DEFERRABLE INITIALLY IMMEDIATE; 
ALTER TABLE "Retroalimentacion" ADD CONSTRAINT "fk_Retroalimentacion_id_Simulacion_Simulacion" FOREIGN KEY ("id_Simulacion") REFERENCES "Simulacion" ("id_Simulacion") ON DELETE NO ACTION ON UPDATE NO ACTION DEFERRABLE INITIALLY IMMEDIATE; 
ALTER TABLE "Retroalimentacion" ADD CONSTRAINT "fk_Retroalimentacion_id_ComportamientoVial_ComportamientoVial" FOREIGN KEY ("id_ComportamientoVial") REFERENCES "ComportamientoVial" ("id_ComportamientoVial") ON DELETE NO ACTION ON UPDATE NO ACTION DEFERRABLE INITIALLY IMMEDIATE; 
ALTER TABLE "MetricaDesempeno" ADD CONSTRAINT "fk_MetricaDesempeno_id_Simulacion_Simulacion" FOREIGN KEY ("id_Simulacion") REFERENCES "Simulacion" ("id_Simulacion") ON DELETE NO ACTION ON UPDATE NO ACTION DEFERRABLE INITIALLY IMMEDIATE; 
ALTER TABLE "Reporte" ADD CONSTRAINT "fk_Reporte_id_Simulacion_Simulacion" FOREIGN KEY ("id_Simulacion") REFERENCES "Simulacion" ("id_Simulacion") ON DELETE NO ACTION ON UPDATE NO ACTION DEFERRABLE INITIALLY IMMEDIATE; 
ALTER TABLE "Recompensa" ADD CONSTRAINT "fk_Recompensa_id_Usuario_Usuario" FOREIGN KEY ("id_Usuario") REFERENCES "Usuario" ("id_Usuario") ON DELETE NO ACTION ON UPDATE NO ACTION DEFERRABLE INITIALLY IMMEDIATE; 
ALTER TABLE "ReglaTransito" ADD CONSTRAINT "fk_ReglaTransito_id_Escenario_Escenario" FOREIGN KEY ("id_Escenario") REFERENCES "Escenario" ("id_Escenario") ON DELETE NO ACTION ON UPDATE NO ACTION DEFERRABLE INITIALLY IMMEDIATE; 
ALTER TABLE "EventoVial" ADD CONSTRAINT "fk_EventoVial_id_Simulacion_Simulacion" FOREIGN KEY ("id_Simulacion") REFERENCES "Simulacion" ("id_Simulacion") ON DELETE NO ACTION ON UPDATE NO ACTION DEFERRABLE INITIALLY IMMEDIATE; 
ALTER TABLE "Infraccion" ADD CONSTRAINT "fk_Infraccion_id_Simulacion_Simulacion" FOREIGN KEY ("id_Simulacion") REFERENCES "Simulacion" ("id_Simulacion") ON DELETE NO ACTION ON UPDATE NO ACTION DEFERRABLE INITIALLY IMMEDIATE; 
\end{lstlisting}

\section{Diseño de interfaz: wireframes}
Se han diseñado las maquetas de baja fidelidad para las vistas principales del sistema.

\begin{itemize}
    \item \textbf{Pantalla 1 - Inicio de Sesión / Registro:} Formulario de acceso. Muestra campos para correo, contraseña y enlace de recuperación.
\begin{figure}[H]
    \centering
    \includegraphics[width=0.8\textwidth]{Pantalla 1 - Inicio de Sesión Registro.png}
    \caption{Pantalla 1 - Inicio de Sesión / Registro}
\end{figure}

\item \textbf{Pantalla 2 - Dashboard Principal:} Vista post-login donde el usuario (conductor o instructor) ve el resumen de sus estadísticas y un botón destacado para iniciar nueva práctica.
\begin{figure}[H]
    \centering
    \includegraphics[width=0.8\textwidth]{Pantalla 2 - Dashboard Principal.png}
    \caption{Pantalla 2 - Dashboard Principal}
\end{figure}

\item \textbf{Pantalla 3 - Selección de Escenario (Core):} Interfaz donde el usuario configura la simulación, eligiendo dificultad y clima antes de cargar el entorno virtual.
\begin{figure}[H]
    \centering
    \includegraphics[width=0.8\textwidth]{Pantalla 3 - Selección de Escenario (Core).png}
    \caption{Pantalla 3 - Selección de Escenario (Core)}
\end{figure}

\item \textbf{Pantalla 4 - Reporte de Retroalimentación:} Visualización de puntajes, listado de errores específicos cometidos en la sesión y sugerencias de mejora pedagógicas generadas por la IA.
\begin{figure}[H]
    \centering
    \includegraphics[width=0.8\textwidth]{Pantalla 4 - Reporte de Retroalimentación.png}
    \caption{Pantalla 4 - Reporte de Retroalimentación}
\end{figure}

\end{itemize}

\section{Plan de trabajo: cronograma e hitos}
\subsection{Cronograma semanal}
\begin{table}[H]
\centering
\resizebox{\textwidth}{!}{
\begin{tabular}{|p{1.5cm}|p{6cm}|p{3.5cm}|p{3.5cm}|}
\hline
\textbf{Semana} & \textbf{Tarea principal} & \textbf{Responsable} & \textbf{Entregable} \\ \hline
5 & Entrega 1A: requerimientos, arquitectura C4, MER, diccionario de datos, DDL y wireframes & Todo el equipo & PDF Entrega 1A en LMS \\ \hline
6 & Configuración del entorno, repositorio GitHub, estructura del proyecto y conexión PostgreSQL & Zamora D. & Repositorio inicial \\ \hline
7 & Desarrollo de la base de datos PostgreSQL y script DDL completo & Umaginga J. & PE Unidad II \\ \hline
8 & Tutoría, revisión de avances y corrección de observaciones & Todo el equipo & Avance validado \\ \hline
9 & Evaluación Parcial 1 Individual & Todo el equipo & EP Primer Corte \\ \hline
10 & Implementación del módulo de autenticación, usuarios y roles & Todo el equipo & Avance PFC Módulo Datos \\ \hline
11 & Desarrollo CRUD de Escenarios, Simulaciones y Reglas de Tránsito & Cruz Perez J. & PE Unidad III \\ \hline
12 & Integración Backend–Base de Datos y pruebas funcionales & Umaginga J. & Módulos integrados \\ \hline
13 & Desarrollo del Frontend responsivo e interfaces principales & Zamora D. & Exposición Back-End \\ \hline
14 & Implementación de API REST y módulo de análisis con IA & Todo el equipo & Sistema integrado \\ \hline
15 & Entrega 1B: implementación completa y documentación & Todo el equipo & Entrega 1B \\ \hline
16 & Pruebas de aceptación, corrección de errores y optimización & Todo el equipo & PE Unidad IV \\ \hline
17 & Defensa oral del Proyecto Final de Curso & Todo el equipo & EP Segundo Corte + Defensa \\ \hline
\end{tabular}}
\end{table}

\subsection{Asignación de roles del equipo}
\begin{table}[H]
\centering
\begin{tabular}{|p{4cm}|p{3cm}|p{7cm}|}
\hline
\textbf{Integrante} & \textbf{Rol} & \textbf{Responsabilidades} \\ \hline
Zamora Bumbila Diego A. & Líder Técnico & Coordinación del proyecto, arquitectura C4, integración de módulos, revisión de código, control de versiones y documentación técnica. \\ \hline
Umaginga Arevalo Jefferson M. & Dev Backend y Base de Datos & PostgreSQL, script DDL, modelos, controladores, autenticación, API REST y lógica de negocio. \\ \hline
Cruz Perez Justyn K. & Dev Frontend & Diseño de interfaces, wireframes, HTML5, CSS3, JavaScript, Bootstrap, formularios y experiencia de usuario. \\ \hline
\end{tabular}
\end{table}

\section{Estructura del repositorio Git}
El proyecto contará con un repositorio alojado en GitHub. La URL se encuentra en la portada de este documento.

Estructura requerida al momento de la entrega:
\begin{lstlisting}[language=bash]
README.md                      # Descripcion del proyecto e instrucciones
docs/entrega-1a.pdf            # PDF de este informe
docs/diagramas/                # PNG de diagramas C4, MER, Wireframes
docs/adr/ADR-001-tecnologia.md # ADR de la decision del backend
database/schema.sql            # Script DDL de la base de datos PostgreSQL
.gitignore                     # Archivos ignorados (vendor/, node_modules/)
.env.example                   # Nombres de variables de entorno (sin valores)
\end{lstlisting}
\newpage
\section{Referencias bibliográficas}
\begin{thebibliography}{99}

\bibitem{sommerville}
Sommerville, I. (2015). \textit{Ingeniería del software} (10ma ed.). Pearson Education.

\bibitem{brown}
Brown, S. (2023). \textit{The C4 model for visualising software architecture}. Leanpub. https://c4model.com/

\bibitem{shirole2025}
Shirole, V., Shahade, A. K., \& Deshmukh, P. V. (2025). A comprehensive review on data-driven driver behaviour scoring in vehicles: technologies, challenges and future directions. \textit{Discover Artificial Intelligence}, 5(1), 26. \url{https://doi.org/10.1007/S44163-025-00244-6}

\bibitem{afridi2025}
Afridi, A. S., Kafy, A., Moon, M. N. N., \& Shakil, M. S. (2025). A multi-class driver behavior dataset for real-time detection and road safety enhancement. \textit{Data Brief}, 60, 111529. \url{https://doi.org/10.1016/J.DIB.2025.111529}

\bibitem{nixon2021}
Nixon, R. (2021). \textit{Learning PHP, MySQL \& JavaScript} (6th ed.). O'Reilly Media.


\end{thebibliography}


\end{document}

